\documentclass{article}

\usepackage{iclr2027_conference,times}
\usepackage{amsmath,amssymb,amsthm,mathtools}
\usepackage{xcolor}
\usepackage{mathrsfs}
\usepackage{booktabs}
\usepackage{enumitem}
\usepackage{microtype}
\usepackage[hidelinks]{hyperref}
\usepackage{url}
\usepackage[nameinlink,noabbrev]{cleveref}

\newtheorem{definition}{Definition}[section]
\newtheorem{theorem}[definition]{Theorem}
\newtheorem{proposition}[definition]{Proposition}
\newtheorem{lemma}[definition]{Lemma}
\newtheorem{corollary}[definition]{Corollary}
\newtheorem{example}[definition]{Example}
\newtheorem{remark}[definition]{Remark}

\crefname{definition}{definition}{definitions}
\crefname{theorem}{theorem}{theorems}
\crefname{proposition}{proposition}{propositions}
\crefname{lemma}{lemma}{lemmas}
\crefname{corollary}{corollary}{corollaries}
\crefname{example}{example}{examples}
\Crefname{definition}{Definition}{Definitions}
\Crefname{theorem}{Theorem}{Theorems}
\Crefname{proposition}{Proposition}{Propositions}
\Crefname{lemma}{Lemma}{Lemmas}
\Crefname{corollary}{Corollary}{Corollaries}
\Crefname{example}{Example}{Examples}

\newcommand{\Id}{\operatorname{id}}
\newcommand{\Eq}{\operatorname{Eq}}
\newcommand{\Hit}{\operatorname{Hit}}
\newcommand{\Max}{\operatorname{Max}}
\newcommand{\Min}{\operatorname{Min}}
\newcommand{\Opt}{\operatorname{Opt}}
\newcommand{\supp}{\operatorname{supp}}
\newcommand{\TV}{\operatorname{TV}}
\newcommand{\Set}{\mathbf{Set}}
\newcommand{\SBor}{\mathbf{SBor}}
\newcommand{\Two}{\mathbf{2}}
\newcommand{\Quot}{\mathbf{Quot}}
\newcommand{\Supp}{\mathbf{Supp}}
\newcommand{\Up}{\mathbf{Up}}
\newcommand{\Bool}{\mathbb{B}}
\newcommand{\Prob}{\mathsf{Prob}}
\newcommand{\State}{\mathsf{State}}
\newcommand{\Dec}{\mathsf{Dec}}
\newcommand{\Mech}{\mathsf{Mech}}
\newcommand{\Loc}{\mathsf{Loc}}
\newcommand{\Defect}{\mathsf{D}}
\newcommand{\NN}{\mathsf{N}}
\newcommand{\Rcal}{\mathcal{R}}
\newcommand{\Mcal}{\mathcal{M}}
\newcommand{\Hcal}{\mathcal{H}}
\newcommand{\Ccal}{\mathcal{C}}
\newcommand{\Bcal}{\mathcal{B}}
\newcommand{\Gcal}{\mathcal{G}}

\title{Mechanistic Localization Is\\
Distribution-Relative:\\
Receiver Sufficiency, Context Invariance,\\
and the Support Shadow}
\author{Anonymous authors}

\begin{document}
\maketitle

\begin{abstract}
Mechanistic localization is distribution-relative. A component set does not determine the optimal rule obtainable from the internal representation it exposes: even with the network, receiver, phenomenon, and literal support fixed, changing only probability weights can reverse the Bayes-optimal rule. We therefore include the realized probability state in the localization problem and distinguish a structural location $K$, a receiver-level presentation $(K,h)$---the selected receiver set together with a decoder---and the induced action $[hQ_K]_\mu$. The receiver map determines what state is exposed; an admissible decoder class determines how that state may be used. We take all measurable receiver-level decoders as the maximal receiver-sufficiency class and recover architecture-specific continuation classes by restriction. This separation yields two different laws. Passive receiver refinement can only improve predictive adequacy, but it does not order context invariance. Under log loss, the value of revealing context is $I(Y_\Phi;C\mid Q_K(X))$, and adding receiver $j$ changes this value by $I(Y_\Phi;Q_j\mid Q_K,C)-I(Y_\Phi;Q_j\mid Q_K)$; either sign occurs. Exact localization is strictly coarser. In finite $0$--$1$ settings, zero risk forgets probability weights and reduces to support-level factorization, functional dependency, and decision-reduct structure. Fixing the network does not fix this exact localization family: one Boolean OR network with fixed coordinate receivers realizes every nonempty upward-closed exact localization family by varying support alone. A circuit label is therefore not a complete mechanistic localization claim without the realized distribution, access, admissible downstream use, phenomenon, and adequacy criterion under which it is asserted.
\end{abstract}

\section{Introduction}
\label{sec:introduction}

A circuit-localization claim names a subset of internal components and asks whether the exposed state through those components suffices for a phenomenon. Decision-theoretic representation theory asks whether an exposed representation is Bayes-sufficient for a specified distribution and loss \citep{sevetlidis2026bayes}. That fixes the probability law while asking what the representation makes available. Mechanistic localization then asks what survives when the network, receiver, phenomenon, loss, and literal support stay fixed while the realized probability law changes.

Even the smallest finite example shows that the receiver-level rule need not survive. Let
\[
X=\{0,1\}^2,\qquad
\Phi(x_1,x_2)=x_1,\qquad
Q_2(x_1,x_2)=x_2.
\]
For $0<\eta<1/2$, let $\mu_\eta$ put mass $(1-\eta)/2$ on each diagonal point $00,11$ and mass $\eta/2$ on each off-diagonal point. Let $\nu_\eta$ reverse those weights. The represented state space, receiver, phenomenon, and literal support are unchanged:
\[
\supp\mu_\eta=\supp\nu_\eta=X.
\]
Under $0$--$1$ loss, however, the optimal rule through the same receiver reverses,
\[
h_{\mu_\eta}(z)=z,\qquad h_{\nu_\eta}(z)=1-z,
\]
and has error $\eta$ in either state. Their equal mixture is uniform; there $Q_2$ carries no predictive information about $\Phi$ and the minimum error is $1/2$. Passing from the probability law to its support erases exactly the information that selects the receiver-level rule. Thus neither a component set nor a support-relative circuit determines the receiver-level realization selected by the probability distribution. We therefore treat that distribution as part of the realized computational state.

Recent results often grouped under ``circuit variability'' vary different objects. Circuit structure changes under input variation, resampling, and prompt rephrasing, while transfer and interchange tests can make some structural differences functionally weak \citep{bayat2026many,wu2026variance}. A behavior can admit several sparse or otherwise admissible circuit descriptions \citep{chen2026allcircuits,meloux2025identifiable}; formal verification can separately certify robustness or minimality for a declared candidate \citep{hadad2026formal}. Mediator choice changes the exposed internal state, and interventions can move that state away from the natural distribution without changing learned parameters \citep{mueller2026mediator,grant2026divergent}. These are not one kind of instability: they change different arguments of the mechanistic claim.

Existing formalisms already expose some of these arguments. Causal abstraction fixes maps between computational systems; verification-based discovery certifies a declared circuit on an explicit domain; MIB fixes benchmark tracks and evaluation targets for circuit and causal-variable localization \citep{geiger2025causal,hadad2026formal,mueller2025mib}. Here we hold the ambient network, receiver access, phenomenon, admissible receiver-to-action class, and decision criterion fixed and vary the realized probability state. This isolates a precise question: which parts of a localization claim survive a change in the law governing the realized states?

A deterministic network supplies an ambient process $X^\theta$. A source law $\mu_0$ induces a state-decorated process $\widetilde X^{\theta,\mu_0}$. At a represented cut, the analysis declares receiver maps $Q_j$ and a phenomenon $\Phi$. Operationally, $Q_K(X)$ is the internal representation made available to the analysis, while a decoder $h$ specifies how that representation is used to predict or realize the target phenomenon. The exposed state and its admissible downstream use are separate choices. For receiver set $K$, let $\Dec_{\mathsf A}(K)$ be the measurable maps from $Z_K$ to an action space $\mathsf A$, and let $\mathscr H_K\subseteq\Dec_{\mathsf A}(K)$ denote the admissible decoder hypothesis class induced by a declared downstream architecture or continuation protocol. The maximal choice $\mathscr H_K=\Dec_{\mathsf A}(K)$ isolates what is realizable from the receiver state before any further downstream architectural restriction. The main theory works at this maximal receiver-sufficiency level; architecture-specific localization is obtained by restricting $\mathscr H$.

For $h\in\mathscr H_K$, the resulting hierarchy is
\begin{equation}
\boxed{
K\longrightarrow (K,h)\longrightarrow [hQ_K]_\mu .
}
\label{eq:intro-hierarchy}
\end{equation}
In ordinary terms, these are the selected internal variables, those variables together with a decoder, and the predictor that this decoder induces on the realized distribution. Location is not presentation, and presentation is not realized action. The first object records where the analysis reads; the second records one admissible rule applied to that state; the third records the induced action on the realized population up to almost-sure equality. The opening example holds $K$ fixed while changing the optimal presentation and realized action. A component set is therefore not a complete localization specification.

This distinction yields two different laws. Inside one fixed passive access family,
\[
K\subseteq K'
\quad\Longrightarrow\quad
\delta_\mu(K')\le\delta_\mu(K):
\]
more exposed receiver state cannot worsen optimal prediction. Context-invariant realization obeys no corresponding monotonicity law. If $C$ labels contexts, then under log loss
\[
\Gamma_K^{\log}=I(Y_\Phi;C\mid Q_K(X)),
\]
and adding receiver $j$ changes this quantity by
\begin{equation}
\boxed{
\Gamma_{K\cup\{j\}}^{\log}-\Gamma_K^{\log}
=
I(Y_\Phi;Q_j\mid Q_K,C)
-
I(Y_\Phi;Q_j\mid Q_K).
}
\label{eq:intro-increment}
\end{equation}
Either sign occurs. Receiver refinement therefore orders predictive adequacy but does not order context invariance: more internal information can make a common rule easier or harder to sustain across distributions.

Exact localization is a strictly coarser problem. Almost-sure exact realization depends only on measure class; in finite $0$--$1$ problems, zero risk depends only on support and becomes support-level factorization, functional dependency, and decision-reduct structure. This coarse-graining destroys probability-weight information by construction. It also destroys any hope that architecture alone selects a canonical exact circuit at the receiver-sufficiency level. We prove that one fixed Boolean OR network, with fixed coordinate receivers and a fixed Boolean target, realizes every nonempty upward-closed exact localization family by varying support alone. The paper therefore establishes both sides of the distribution-relative picture: quantitative receiver adequacy has a monotone refinement law, while context invariance and exact circuit identity remain distribution-relative. The appendices carry the architecture-specialization result, longer finite geometry, categorical details, and exact descent constructions.

\section{Realized localization and predictive refinement}
\label{sec:realized-access-presentations}
\label{sec:stateful-processes}
\label{sec:mechanism-presentations}

Probability state is not metadata attached to a neural computation. It determines which receiver-level rule is optimal and therefore belongs in the realized localization problem. We include it in the process before defining receiver sufficiency.

\paragraph{Realized process.}
Let $\SBor$ be the category of standard Borel spaces and measurable deterministic maps, and let
\[
X^\theta:[L]\to\SBor,\qquad [L]=\{0<1<\cdots<L\},
\]
be a deterministic feed-forward neural process. Write $A_q^\theta:X_0\to X_q$ for the prefix to cut $q$. A source law $\mu_0\in\Prob(X_0)$ induces $\mu_q=(A_q^\theta)_*\mu_0$. Thus one input distribution induces a distribution over represented states at every cut of the same fixed network. The following lift packages the network maps and these transported distributions into a single realized process. The probability functor lifts the ambient network to
\begin{equation}
\widetilde X^{\theta,\mu_0}:[L]\to\int\Prob,\qquad
q\mapsto(X_q,\mu_q),
\label{eq:stateful-lift}
\end{equation}
where $f:(X,\mu)\to(Y,\nu)$ in $\int\Prob$ means that $f$ is measurable and $f_*\mu=\nu$. Functoriality of pushforward makes every network continuation state preserving in this lifted sense. Appendix~A gives the formal lift proposition and proof \citep{fritz2020synthetic}.

Equation~\eqref{eq:stateful-lift} describes source-law variation through one fixed network. An internal intervention can instead supply a new boundary law at a represented cut; after that state is supplied, the same downstream construction applies to the rebased suffix. This paper studies the resulting realized process. Modeling the intervention operation that connects the natural and altered processes is a different comparison problem.

\paragraph{Declared access and presentations.}
Localization is relative to an access interface. Fix a represented cut---an internal layer or interface at which the analysis reads the computation---and suppress its index. A mechanistic protocol declares a finite receiver family of measurable maps
\[
Q_j:X\to Z_j,\qquad j\in J,
\]
into standard Borel receiver spaces. A receiver is any declared internal variable or readout used by the localization analysis; neurons, attention heads, sparse features, edges, and causal mediators instantiate different choices of $Q_j$ \citep{mueller2025mib,mueller2026mediator}. The protocol also declares a target phenomenon $\Phi:X\to Y$ into a standard Borel target space $Y$, representing the behavior or quantity to be localized. For $K\subseteq J$, set
\[
Z_K=\prod_{j\in K}Z_j,\qquad Q_K=\langle Q_j\rangle_{j\in K}.
\]
If $K\subseteq K'$, coordinate projection $p_{K'K}:Z_{K'}\to Z_K$ satisfies
\begin{equation}
Q_K=p_{K'K}Q_{K'}.
\label{eq:receiver-garbling}
\end{equation}
These choices define different receiver interfaces rather than interchangeable descriptions of one fixed access map.

Fix a standard Borel action space $\mathsf A$ and measurable loss $\ell:Y\times\mathsf A\to[0,\infty]$. Define the maximal receiver-level decoder class
\[
\Dec_{\mathsf A}(K):=\{h:Z_K\to\mathsf A\mid h\text{ measurable}\}.
\]
A concrete downstream architecture or continuation protocol can instead declare a restricted decoder hypothesis class $\mathscr H_K\subseteq\Dec_{\mathsf A}(K)$. These are distinct sources of limitation: $Q_K$ determines what information is available at the receiver interface, while $\mathscr H_K$ determines which functions of that information the downstream system can realize. The main theory takes $\mathscr H_K=\Dec_{\mathsf A}(K)$ to isolate receiver sufficiency itself. Architecture-specific localization is obtained by restricting $\mathscr H$.

A receiver-level presentation is a receiver set together with one admissible decoder, $(K,h)$ with $h\in\mathscr H_K$. Two presentations are state-relative equivalent when
\[
(K,h)\equiv_\mu(K',h')
\quad\Longleftrightarrow\quad
hQ_K=h'Q_{K'}\quad\mu\text{-a.s.}
\]
Thus presentations are identified when they induce the same predictor on the realized population, even if they use different internal variables. This is the formal version of \cref{eq:intro-hierarchy}. Different locations can induce the same realized action: on $X=\{0,1\}$, the accesses $Q_1(x)=x$ and $Q_2(x)=1-x$ paired with $h_1(z)=z$ and $h_2(z)=1-z$ satisfy $h_1Q_1=h_2Q_2$ pointwise. The opening example gives the converse: one fixed location supports different optimal rules under different states. Location, presentation, and realized action are therefore distinct levels of mechanistic description. Appendix~A records the presentation category, architecture-specific decoder subfamilies, and measure-class properties.

\paragraph{Receiver Bayes risk.}
For a declared decoder family $\mathscr H$, define
\begin{equation}
\delta_{\mu,\mathscr H}(K)
:=
\inf_{h\in\mathscr H_K}
\mathbb E_\mu\!\left[\ell\bigl(\Phi(X),h(Q_K(X))\bigr)\right].
\label{eq:receiver-risk-general}
\end{equation}
The maximal receiver-sufficiency risk used throughout the main results is
\begin{equation}
\delta_\mu(K)
:=
\delta_{\mu,\Dec_{\mathsf A}}(K)
=
\inf_{h:Z_K\to\mathsf A}
\mathbb E_\mu\!\left[\ell\bigl(\Phi(X),h(Q_K(X))\bigr)\right].
\label{eq:receiver-risk}
\end{equation}
For fixed $\mu$, this is exactly the Bayes risk of the representation exposed by $K$; Bayes-sufficient representation theory studies the corresponding fixed-law decision problem \citep{sevetlidis2026bayes}. Here $K$ ranges over a declared internal receiver family and the paper compares the resulting risk profile across realized probability states. For a tolerance $\varepsilon\ge0$, write
\[
\Rcal_{\mu,\varepsilon}:=\{K\subseteq J:\delta_\mu(K)\le\varepsilon\}.
\]
The full risk profile and the thresholded circuit family are different objects.

Here refinement is passive: $K'$ exposes additional receiver variables from the same realized computation rather than modifying the computation itself. Equation~\eqref{eq:receiver-garbling} exhibits the smaller receiver state as a deterministic garbling of the larger one, so receiver inclusion is a Blackwell comparison of experiments \citep{blackwell1951comparison,fritz2023experiments}.
\begin{proposition}[Receiver inclusion is Blackwell refinement]
\label{thm:receiver-defect-monotonicity}
If $K\subseteq K'$, then
\[
\delta_\mu(K')\le \delta_\mu(K).
\]
\end{proposition}
\begin{proof}
For any rule $h:Z_K\to\mathsf A$, let $h'=h\circ p_{K'K}$. Equation~\eqref{eq:receiver-garbling} gives $h'Q_{K'}=hQ_K$ pointwise, so both rules have the same risk. Taking infima gives the claim.
\end{proof}

The proof uses only closure under forgetting added receiver coordinates. Hence every architecture-specific family satisfying
\[
h\in\mathscr H_K\quad\Longrightarrow\quad h\circ p_{K'K}\in\mathscr H_{K'}
\qquad(K\subseteq K')
\]
inherits the same monotonicity law for $\delta_{\mu,\mathscr H}$. In ML terms, the decoder class available at $K'$ must be able to ignore the newly exposed coordinates. Appendix~A states this as a subfunctor result. Passive receiver access therefore has a monotonicity law whenever the admissible downstream class respects receiver refinement. If $K$ instead specifies an ablation, patching, retention, or rerun protocol, changing $K$ can change the state passed to the continuation; structural inclusion alone then does not imply Blackwell refinement.

Distribution relativity does not make receiver adequacy unstable. Let
\[
d_{\TV}(\mu,\nu):=\sup_A|\mu(A)-\nu(A)|.
\]
If $0\le\ell\le1$, then
\begin{equation}
\sup_{K\subseteq J}|\delta_\mu(K)-\delta_\nu(K)|
\le d_{\TV}(\mu,\nu).
\label{eq:tv-main}
\end{equation}
Indeed, for each fixed rule the pointwise loss lies in $[0,1]$, so its expectations under $\mu$ and $\nu$ differ by at most $d_{\TV}(\mu,\nu)$; infimizing and exchanging $\mu,\nu$ gives \eqref{eq:tv-main}. Thus the entire receiver-risk profile is $1$-Lipschitz in total variation, even though optimizer identities and thresholded circuit families can change discontinuously when a decision boundary is crossed. Appendix~B gives the corresponding threshold-stability and finite polyhedral selection geometry.

\section{Context-invariant realization}
\label{sec:contextuality}

Predictive sufficiency and context-invariant realization are different properties. Receiver refinement monotonically improves the first. It need not improve the second. Continue with the maximal receiver decoder class from \cref{sec:realized-access-presentations}; fix $X$, the receiver family, $\Phi$, the action space $\mathsf A$, and the loss $\ell$, and vary only the probability law.

Let $C$ label experimental contexts---such as input populations, domains, prompt families, or other distributional conditions---with $\Pr(C=c)=\lambda_c>0$ and $X\mid C=c\sim\mu_c$. Write
\[
\bar\mu=\sum_c\lambda_c\mu_c.
\]
When $C$ is observed, each context may use its own optimal rule through $Q_K$. When $C$ is hidden, one rule must act on the pooled law. We measure the cost of requiring one receiver-level decoder to serve all contexts by
\begin{equation}
\Gamma_K
:=
\delta_{\bar\mu}(K)-\sum_c\lambda_c\delta_{\mu_c}(K)
\ge0.
\label{eq:context-value}
\end{equation}
It is the excess Bayes risk incurred when the context label is hidden.

\begin{proposition}[Context invariance is a common-optimum condition]
\label{thm:shared-decoder-criterion}
Assume the context and pooled minima are attained. Then
\[
\Gamma_K=0
\quad\Longleftrightarrow\quad
\bigcap_c\Opt_{\mu_c}(K)\neq\varnothing,
\]
where $\Opt_{\mu_c}(K)$ is the set of risk-minimizing rules through $Q_K$.
\end{proposition}
\begin{proof}
A rule optimal in every context attains equality in \cref{eq:context-value}. Conversely, if equality holds and $h^*$ minimizes pooled risk, then
\[
\sum_c\lambda_c R_{\mu_c}(K,h^*)
=
\sum_c\lambda_c\delta_{\mu_c}(K).
\]
Every term on the left is at least the corresponding optimum, and every $\lambda_c$ is positive. Hence $h^*$ is optimal in every context.
\end{proof}

Nothing in this argument depends on maximality of the decoder class. For any fixed architecture-specific family $\mathscr H_K$, replacing $\delta_\mu$ by $\delta_{\mu,\mathscr H}$ gives the same common-optimum criterion with optimizers taken inside $\mathscr H_K$. The criterion therefore isolates the internal invariance question at either level: once the receiver map and admissible downstream class are fixed, does one decoder remain optimal across realized states? This differs from invariant prediction and invariant risk minimization, which impose related common-predictor conditions in different problem settings \citep{arjovsky2019irm,ruan2022covariate}.

Log loss makes the maximal receiver-sufficiency penalty exact. Assume $Y_\Phi$ and $C$ are finite, take $\mathsf A=\Prob(Y)$, and use
\[
\ell_{\log}(y,r)=-\log r(y).
\]
Because the maximal class contains the Bayes action, the optimal rule is the conditional law of $Y_\Phi$, so the pooled and context-aware risks are conditional entropies. Since log loss is unbounded, the bounded-loss TV statement from \cref{sec:realized-access-presentations} is not invoked here.

\begin{theorem}[Context value and receiver increment under log loss]
\label{thm:contextuality-cmi}
\label{thm:contextuality-receiver-increment}
Assume the displayed conditional entropies and mutual informations are finite. Then
\begin{equation}
\Gamma_K^{\log}
=
I(Y_\Phi;C\mid Q_K(X)).
\label{eq:context-cmi}
\end{equation}
For $j\notin K$, writing $Z=Q_K(X)$ and $W=Q_j(X)$,
\begin{equation}
\boxed{
\Gamma_{K\cup\{j\}}^{\log}-\Gamma_K^{\log}
=
I(Y_\Phi;W\mid Z,C)
-
I(Y_\Phi;W\mid Z).
}
\label{eq:receiver-increment}
\end{equation}
The increment can be positive or negative.
\end{theorem}
\begin{proof}
For log loss,
\[
\delta_{\bar\mu}^{\log}(K)=H(Y_\Phi\mid Z),
\qquad
\sum_c\lambda_c\delta_{\mu_c}^{\log}(K)=H(Y_\Phi\mid Z,C),
\]
which gives \cref{eq:context-cmi}. Also,
\[
\Gamma_{K\cup\{j\}}^{\log}-\Gamma_K^{\log}
=
I(Y_\Phi;C\mid Z,W)-I(Y_\Phi;C\mid Z).
\]
Expanding $I(Y_\Phi;C,W\mid Z)$ in the two chain-rule orders gives \cref{eq:receiver-increment}.
\end{proof}

Equation~\eqref{eq:context-cmi} gives a direct interpretation: $\Gamma_K^{\log}$ is exactly the target-relevant information carried by context that remains after the receiver representation is known. Equation~\eqref{eq:receiver-increment} then rules out any general monotonicity of context invariance under receiver refinement. More internal information can make a common rule easier to sustain or harder to sustain across distributions.

Both failures occur in elementary binary systems. For an increase, let $C$ and $Y_\Phi$ be independent fair bits, let $N\sim\mathrm{Bernoulli}(p)$ be independent with $0<p<1/2$, and set
\[
Q_1=Y_\Phi\oplus N,\qquad Q_2=Y_\Phi\oplus C.
\]
With $K=\{1\}$, $C$ is independent of $(Y_\Phi,Q_1)$, so $\Gamma_{\{1\}}^{\log}=0$. Marginally, $Q_2$ is also independent of $(Y_\Phi,Q_1)$; pooled predictive risk does not improve. After conditioning on $Q_2$, however, $C=Y_\Phi\oplus Q_2$, and
\[
\Gamma_{\{1,2\}}^{\log}
=I(Y_\Phi;C\mid Q_1,Q_2)
=H_2(p)>0.
\]
The added receiver has made the optimal decoder context-dependent.

For a decrease, let $C=Y_\Phi$ be nondegenerate, begin from a constant receiver, and add $Q_j=Y_\Phi$. Then
\[
\Gamma_\varnothing^{\log}=I(Y_\Phi;C)=H(Y_\Phi)>0,
\qquad
\Gamma_{\{j\}}^{\log}=I(Y_\Phi;C\mid Y_\Phi)=0.
\]
Here the added receiver removes the common-rule penalty completely.

The opening two-bit construction gives the same separation under $0$--$1$ loss. The context-specific rules are $z$ and $1-z$, their equal mixture has error $1/2$, and
\[
\Gamma_{\{2\}}=\frac12-\eta>0
\]
even though the two context laws have identical support. Exact support localization therefore cannot see the distinction. The next section identifies precisely what remains after the probability weights are discarded.

\section{The exact support shadow and fixed-network non-canonicity}
\label{sec:exact-shadow}
\label{sec:support-shadow}
\label{sec:support-universality}

The distributional theory distinguishes full-support laws that favor opposite receiver-level rules. Exact localization cannot. We now take the maximal receiver-sufficiency problem to zero loss, isolating what the access itself preserves before any further restriction on downstream implementation. Exact localization is obtained after the realized probability state is coarse-grained to its null-set structure:
\begin{equation}
\boxed{
\text{probability state}
\longrightarrow
\text{measure class}
\underset{\text{finite discrete}}{\cong}
\text{support}
\longrightarrow
\text{functional dependency / reduct}.
}
\label{eq:coarse-chain}
\end{equation}
The first arrow forgets probability weights while retaining null sets; in finite discrete spaces those null sets are equivalent to support.

A receiver family $K$ exactly realizes $\Phi$ under $\mu$ if some measurable $h:Z_K\to Y$ satisfies $\Phi=hQ_K$ $\mu$-almost surely. Write
\[
\Rcal_\Phi^{\mathrm{a.s.}}(\mu)
:=
\{K\subseteq J:\exists h\text{ with }\Phi=hQ_K\ \mu\text{-a.s.}\}.
\]

\begin{proposition}[Exact localization depends only on measure class]
\label{thm:measure-class-invariance}
If $\mu\sim\nu$ have the same null measurable sets, then
\[
\Rcal_\Phi^{\mathrm{a.s.}}(\mu)
=
\Rcal_\Phi^{\mathrm{a.s.}}(\nu).
\]
\end{proposition}
\begin{proof}
For fixed $K$ and $h$, the disagreement set $\{x:\Phi(x)\neq h(Q_K(x))\}$ is null under $\mu$ exactly when it is null under $\nu$.
\end{proof}

Now assume $X$, $Y$, and the realized receiver images are finite, take $\mathsf A=Y$, use $0$--$1$ loss, and let $S=\supp\mu$. Define
\[
\Rcal_\Phi(S)
:=
\{K\subseteq J:\exists h:Q_K(S)\to Y\text{ with }\Phi|_S=hQ_K|_S\},
\]
and for $f:S\to T$ write $\Eq(f)=\{(x,x')\in S^2:f(x)=f(x')\}$. We call this exact support-level object the support shadow because it records which target distinctions survive on the realized support while forgetting how probability mass is distributed across that support.

\begin{proposition}[Finite exact localization is the support shadow]
\label{thm:zero-defect-support-shadow}
\label{thm:quotient-realization-shadow}
\label{prop:discernibility-representation}
For every $K\subseteq J$,
\begin{equation}
\delta_\mu(K)=0
\iff
\Phi|_S\text{ factors through }Q_K|_S
\iff
\Eq(Q_K|_S)\subseteq\Eq(\Phi|_S).
\label{eq:zero-kernel-shadow}
\end{equation}
For each pair with $\Phi(x)\neq\Phi(x')$, let
\[
E(x,x'):=\{j\in J:Q_j(x)\neq Q_j(x')\},
\qquad
\Hcal_\Phi(S):=\{E(x,x'):\Phi(x)\neq\Phi(x')\}.
\]
Each $E(x,x')$ is exactly the set of receiver coordinates capable of distinguishing one target-relevant pair. Then
\begin{equation}
\Rcal_\Phi(S)=\Hit(\Hcal_\Phi(S)),
\label{eq:hitting-shadow}
\end{equation}
where $\Hit$ denotes the family of hitting sets.
\end{proposition}
\begin{proof}
Zero $0$--$1$ risk on finite support is equality at every point of positive mass. A decoder through $Q_K$ exists exactly when $Q_K$ never identifies two support points that $\Phi$ distinguishes. Equivalently, at least one receiver in $K$ separates every phenomenon-distinguished pair.
\end{proof}

Equation~\eqref{eq:zero-kernel-shadow} is the exact access test: if it fails, no downstream architecture receiving only $Q_K(X)$ can realize $\Phi$ on $S$. A restricted decoder class can impose an additional implementation constraint even when it holds. Accordingly, the architecture-specific exact family
\[
\Rcal^{\mathscr H}_\Phi(S)
:=
\{K\subseteq J:\exists h\in\mathscr H_K\text{ with }\Phi|_S=hQ_K|_S\}
\]
satisfies $\Rcal^{\mathscr H}_\Phi(S)\subseteq\Rcal_\Phi(S)$.

The finite zero-loss coarse-graining lands exactly on functional-dependency and decision-reduct structure: equation~\eqref{eq:zero-kernel-shadow} is the database functional dependency $K\to d$ when the receivers are condition attributes and $\Phi$ is the decision attribute, and the inclusion-minimal realizing sets are the corresponding decision reducts / minimal transversals \citep{skowron1992discernibility,beeri1984armstrong,demetrovics1979keys}. This identifies classical reduct structure as the exact support shadow of the distribution-relative localization problem.

Support restriction is antitone. If $S\subseteq T$, then
\begin{equation}
\Rcal_\Phi(T)\subseteq\Rcal_\Phi(S),
\label{eq:support-restriction-main}
\end{equation}
because a decoder valid on $T$ remains valid on $S$. The inclusion-minimal family
\[
\Mcal_\Phi(S):=\Min_{\subseteq}\Rcal_\Phi(S)
\]
is generally an antichain rather than a singleton. Minimality therefore provides no canonical exact circuit. Appendix~C records the blocker formulation, support functoriality, and the finite empirical $g_3$ correspondence.

Fixing the network does not fix the exact localization family. The question is how much variation remains once the architecture, target, and receiver interface are all frozen. Within the upward-closure constraint already forced by passive receiver inclusion, one fixed elementary network realizes every admissible family by support variation alone. Let $J\neq\varnothing$ and fix
\begin{equation}
\NN_J:
\{0,1\}^J
\xrightarrow{\Id}
\{0,1\}^J
\xrightarrow{\mathrm{OR}}
\{0,1\},
\label{eq:fixed-or}
\end{equation}
with coordinate receivers $Q_j(v)=v_j$ and phenomenon $\Phi(v)=\mathrm{OR}(v)$. Classical candidate-key and reduct constructions already show broad realizability of Sperner systems \citep{demetrovics1979keys,demetrovics2014reducts}. Here the architecture, binary coordinate interface, and target function are all frozen; only the admitted support changes. For a family $\mathcal F\subseteq2^J$, let $\Hit(\mathcal F)$ be the sets intersecting every member of $\mathcal F$. If $F\subseteq2^J$ is nonempty and upward closed, define
\begin{equation}
\sigma_{\mathrm{OR}}(F)
:=
\{0^J\}\cup\{\mathbf1_E:E\in\Hit(F)\}.
\label{eq:or-support}
\end{equation}

\begin{theorem}[Fixed-network support universality]
\label{thm:fixed-or-support-universality}
For every nonempty upward-closed $F\subseteq2^J$,
\begin{equation}
\boxed{
\Rcal_\Phi(\sigma_{\mathrm{OR}}(F))=F.
}
\label{eq:or-universality}
\end{equation}
Consequently every nonempty antichain $\mathcal A\subseteq2^J$ occurs as $\Mcal_\Phi(S)$ for some support $S$ of this one fixed network.
\end{theorem}
\begin{proof}
The phenomenon-distinguished pairs are the anchor $0^J$ against witnesses $\mathbf1_E$, and their receiver-separation sets are exactly $E$. Hence the discernibility hypergraph of $\sigma_{\mathrm{OR}}(F)$ is $\Hit(F)$. For an upward family, $\Hit(\Hit(F))=F$: if $K\notin F$, upward closure implies every member of $F$ intersects $J\setminus K$, so $J\setminus K\in\Hit(F)$ and $K$ fails to hit it. Equation~\eqref{eq:hitting-shadow} gives \eqref{eq:or-universality}. For an antichain $\mathcal A$, take $F=\uparrow\mathcal A$.
\end{proof}

This exhausts the upward-closure constraint imposed by passive receiver inclusion: at the exact receiver-sufficiency level, every nonempty upward-closed localization family occurs in the same fixed network under support variation alone, even though the architecture, target, and receiver granularity remain fixed. Architecture therefore does not determine a canonical exact localization family. Restricting the admissible decoder class defines a finer architecture-specific problem; it does not change the universality result for the information already exposed at the receiver interface. Appendix~D gives stronger order-profile consequences. The theorem is a structural boundary on exact localization; empirical multiplicity in a particular language model can have additional causes \citep{meloux2025identifiable,chen2026allcircuits}.

\section{Implications for mechanistic localization}
\label{sec:discussion}

A component set is not a complete mechanism specification at the receiver-sufficiency level. A localization claim fixes the realized probability state $\mu$, receiver or mediator map $Q$, phenomenon $\Phi$, admissible decoder class $\mathscr H$, and decision criterion $\ell$; a selection threshold $\varepsilon$ and concrete rule $h$ also belong to the claim when used. Empirically, these are the input or prompt distribution, mediator or representation being read, downstream architecture class, target behavior, adequacy criterion, selection threshold, and concrete receiver-level rule. Different experimental choices therefore act on different arguments of the problem:
\begin{table}[t]
\centering
\small
\begin{tabular}{@{}p{0.27\linewidth}p{0.19\linewidth}p{0.45\linewidth}@{}}
\toprule
\textbf{What changes} & \textbf{Argument} & \textbf{What can change} \\
\midrule
input population or prompt distribution & $\mu$ & optimal rule, receiver risk, context value, selected circuit \\
mediator or representation granularity & $Q$ & the localization problem itself \\
downstream architecture or continuation class & decoder class $\mathscr H$ & realizable receiver-level rules, optimal risk, selected localization \\
target behavior & $\Phi$ & which distinctions receivers must preserve \\
loss or adequacy criterion & $\ell$ & optimal action and receiver risk \\
selection threshold & $\varepsilon$ & the discrete reported circuit family \\
passive receiver set & $K$ & predictive risk, monotonically under refinement \\
intervention protocol & realized process/state & the induced experiment; passive Blackwell order need not apply \\
\bottomrule
\end{tabular}
\caption{Different experimental changes act on different arguments of a mechanistic-localization claim.}
\label{tab:localization-arguments}
\end{table}

Different reported component subsets do not by themselves establish different realized mechanisms: studies may differ in $\mu$, $Q$, $\mathscr H$, $\Phi$, $\ell$, $\varepsilon$, or the induced process. Conversely, the same subset does not establish the same receiver-level realization, because the optimal rule and induced action can vary with $\mu$. The hierarchy $K\to(K,h)\to[hQ_K]_\mu$ makes both failures explicit.

The Blackwell order applies when
\[
Q_K=p_{K'K}Q_{K'}
\]
on the same realized state. Ablation, patching, retention, or rerun protocols can instead change the state entering downstream computation, so structural inclusion alone does not define a Blackwell refinement; comparison then requires a garbling or simulation relation between the induced experiments.

The maximal decoder class is the downstream-implementation-independent outer layer of receiver sufficiency. It separates two obstructions. An \emph{access obstruction}, $\delta_\mu(K)>0$, means no measurable downstream rule receiving only $Q_K(X)$ can attain zero risk: the required distinction is already lost at the receiver interface. A \emph{decoder-class obstruction}, $\delta_\mu(K)=0$ but $\delta_{\mu,\mathscr H}(K)>0$, means the distinction reaches the receiver but the selected architecture cannot realize the required map. These are different mechanistic failures and should not be collapsed into one notion of ``circuit insufficiency.''

Causal necessity and intervention-defined mechanisms add another axis. Restricting $\mathscr H$ changes which functions may act on a fixed representation; patching or ablation can change the representation distribution itself. Causal abstraction and formal patching supply comparison structure for that problem \citep{geiger2025causal,hadad2026formal}; cross-network work fixes still another comparison problem \citep{sun2026equivalent}. Receiver granularity, downstream use class, and induced process are therefore separate declared parts of the analysis.

At zero loss, context invariance becomes a literal local-to-global compatibility problem. If contexts $S_i$ each admit an exact decoder through one fixed $Q_K$, a decoder on $\bigcup_i S_i$ exists exactly when the local decoders agree on every receiver-image overlap
\[
Q_K(S_i)\cap Q_K(S_j).
\]
These can exceed the images of $S_i\cap S_j$ because noninjective access can identify previously distinct states. Appendix~E gives the descent formulation. Thus $\Gamma_K$ and the exact global-section obstruction are two regimes of the same question: whether context-specific receiver rules extend to one rule.

Architecture-level receiver formalisms constrain the represented process, exposed distinctions, and admissible downstream use before a probability state and adequacy criterion are chosen \citep{rosariofreytes2026architecture}. Here those interfaces are fixed and the realized state varies. The maximal decoder class isolates the dependence forced by state and access; restricting $\mathscr H$ specializes the same construction to a concrete downstream architecture. Mechanistic localization is therefore a statement about a realized computation under declared access and admissible use, not a property attached to a bare subset of model components.

\label{sec:main-text-end}


\subsection*{AI use statement}
In this work, generative AI tools were used to help develop the theoretical and conceptual framework; formulate and refine mathematical claims; explore, draft, and check proof arguments; design validation methodology; and interpret formal and finite validation results. They were also used for literature search and synthesis, generating and editing code for finite regression tests, and drafting and editing the manuscript. They were not used to generate or clean empirical datasets, conduct qualitative data analysis, assist with translation, or run empirical experiments; those tasks are not applicable to this theoretical study. All AI-assisted mathematical claims were checked against explicit arguments or finite regressions where applicable, literature claims were checked against cited sources, and the authors reviewed the final text and take responsibility for the ideas, claims, proofs, citations, code, and presentation.

\subsection*{Reproducibility statement}
The paper is theoretical. Assumptions and proofs of the two headline results appear in the main text. Auxiliary proofs, finite structural results, exact-shadow constructions, and the finite regression suite are supplied in the appendices and accompanying source materials so that the formal claims can be audited independently.

\bibliography{references}
\bibliographystyle{iclr2027_conference}

\appendix
\section{Formal details for realized processes and receiver presentations}
\label{app:process-presentation-details}

This appendix records formal properties used by the main-text typing of realized neural processes and receiver-level mechanism presentations. It also formalizes architecture-specific decoder classes as restrictions of the maximal receiver-level decoder functor: the main text's decoder hypothesis classes appear here as subfunctors. The main text keeps the constructions that drive the argument; categorical and almost-sure bookkeeping is collected here.

\subsection{State transport and the realized neural process}

Let $\SBor$ denote the category of standard Borel spaces and measurable deterministic maps. For a measurable map $f:X\to Y$, pushforward defines
\[
f_*:\Prob(X)\to\Prob(Y),
\qquad
(f_*\mu)(B)=\mu(f^{-1}(B)).
\]
Thus
\begin{equation}
\Prob:\SBor\longrightarrow\Set
\label{eq:app-probability-functor}
\end{equation}
is a covariant functor.

Fix represented cuts
\[
[L]=\{0<1<\cdots<L\}
\]
regarded as a category, and let
\[
X^\theta:[L]\to\SBor
\]
be a deterministic feed-forward neural process. Write
\[
F_{q\rightsquigarrow r}^\theta:=X^\theta(q\to r),
\qquad
A_q^\theta:=F_{0\rightsquigarrow q}^\theta.
\]
For a source state $\mu_0\in\Prob(X_0)$, set
\[
\mu_q:=(A_q^\theta)_*\mu_0.
\]
The category of elements $\int\Prob$ has objects $(X,\mu)$ and morphisms
\[
f:(X,\mu)\to(Y,\nu)
\]
with $f:X\to Y$ measurable and $f_*\mu=\nu$.

\begin{proposition}[The realized neural process]
\label{thm:app-state-decorated-process}
For every deterministic neural process $X^\theta$ and source state $\mu_0$, the pushed-forward states define a unique lift
\[
\widetilde X^{\theta,\mu_0}:[L]\longrightarrow\int\Prob
\]
whose projection to $\SBor$ is $X^\theta$ and whose value at cut $q$ is $(X_q,\mu_q)$.
\end{proposition}

\begin{proof}
For $q\le r$,
\[
(F_{q\rightsquigarrow r}^\theta)_*\mu_q
=(F_{q\rightsquigarrow r}^\theta)_*(A_q^\theta)_*\mu_0
=(A_r^\theta)_*\mu_0
=\mu_r.
\]
Hence every network arrow is a morphism in $\int\Prob$. The source state determines every downstream state through the unique arrow $0\to q$, so the lift is unique.
\end{proof}

\begin{proposition}[Equivalent states remain equivalent downstream]
\label{prop:app-measure-class-pushforward}
Let $\mu$ and $\nu$ be probability measures on $X$ with the same null measurable sets. For every measurable deterministic map $f:X\to Y$,
\[
f_*\mu\sim f_*\nu.
\]
Consequently, if two source states lie in the same measure class, their induced states lie in the same measure class at every represented cut of a deterministic neural process.
\end{proposition}

\begin{proof}
For every measurable $B\subseteq Y$,
\[
(f_*\mu)(B)=0
\Longleftrightarrow
\mu(f^{-1}(B))=0
\Longleftrightarrow
\nu(f^{-1}(B))=0
\Longleftrightarrow
(f_*\nu)(B)=0.
\]
Apply the result to each prefix map $A_q^\theta$.
\end{proof}

\subsection{Receiver refinement as a presentation category}

Fix a finite declared receiver family $J$ at one represented cut, with measurable maps
\[
Q_j:X\to Z_j.
\]
Write
\[
\Bool_J:=(2^J,\subseteq)
\]
for the Boolean receiver poset regarded as a category. For $K\subseteq J$, let
\[
Z_K:=\prod_{j\in K}Z_j,
\qquad
Q_K:=\langle Q_j\rangle_{j\in K},
\]
and for $K\subseteq K'$ let
\[
p_{K'K}:Z_{K'}\to Z_K
\]
be coordinate projection. Fix a standard Borel action space $\mathsf A$ and define
\[
\Dec_{\mathsf A}(K)
:=\{h:Z_K\to\mathsf A\mid h\text{ measurable}\}.
\]
Receiver refinement pulls a rule on $K$ back to one on $K'$,
\begin{equation}
\Dec_{\mathsf A}(K)\longrightarrow\Dec_{\mathsf A}(K'),
\qquad
h\longmapsto h\circ p_{K'K}.
\label{eq:app-decoder-refinement-map}
\end{equation}
These maps are functorial, so
\[
\Dec_{\mathsf A}:\Bool_J\to\Set
\]
is a functor. Its category of elements $\int\Dec_{\mathsf A}$ has objects $(K,h)$ with $K\subseteq J$ and $h:Z_K\to\mathsf A$.

A refinement arrow preserves the induced action pointwise. If $K\subseteq K'$ and $h'=h\circ p_{K'K}$, then
\begin{equation}
h'Q_{K'}=hQ_K.
\label{eq:app-refinement-preserves-action}
\end{equation}

\subsection{Architecture-specific decoder subfunctors}

The maximal decoder functor separates receiver access from any particular downstream implementation. Categorically, refinement-compatible decoder hypothesis classes are exactly subfunctors of the maximal decoder functor,
\[
\mathscr H\hookrightarrow\Dec_{\mathsf A}.
\]
Concretely, this consists of subsets
\[
\mathscr H_K\subseteq\Dec_{\mathsf A}(K)
\]
for each $K\subseteq J$ such that
\begin{equation}
h\in\mathscr H_K
\quad\Longrightarrow\quad
h\circ p_{K'K}\in\mathscr H_{K'}
\qquad(K\subseteq K').
\label{eq:app-architecture-refinement-closure}
\end{equation}
This is exactly the closure needed for passive receiver refinement to preserve the admissible downstream class.

For a probability state $\mu$ and measurable loss $\ell:Y\times\mathsf A\to[0,\infty]$, define
\[
\delta_{\mu,\mathscr H}(K)
:=
\inf_{h\in\mathscr H_K}
\mathbb E_\mu\!\left[\ell\bigl(\Phi(X),h(Q_K(X))\bigr)\right].
\]

\begin{proposition}[Refinement monotonicity for decoder subfunctors]
\label{prop:app-architecture-refinement-monotonicity}
If $\mathscr H\hookrightarrow\Dec_{\mathsf A}$ is a refinement-compatible decoder subfunctor and $K\subseteq K'$, then
\[
\delta_{\mu,\mathscr H}(K')\le\delta_{\mu,\mathscr H}(K).
\]
\end{proposition}

\begin{proof}
For $h\in\mathscr H_K$, closure \eqref{eq:app-architecture-refinement-closure} gives $h':=h\circ p_{K'K}\in\mathscr H_{K'}$. Equation~\eqref{eq:app-refinement-preserves-action} gives $h'Q_{K'}=hQ_K$ pointwise, so the two rules have the same risk. Taking the infimum over $h\in\mathscr H_K$ gives the claim.
\end{proof}

The maximal case is $\mathscr H=\Dec_{\mathsf A}$. Smaller subfunctors specialize the receiver-sufficiency problem to restricted downstream computation classes without changing the access maps themselves. A family that is not closed under \eqref{eq:app-architecture-refinement-closure} need not inherit the refinement monotonicity law.

In the finite exact setting, define
\[
\Rcal^{\mathscr H}_\Phi(S)
:=
\{K\subseteq J:\exists h\in\mathscr H_K\text{ with }\Phi|_S=hQ_K|_S\}.
\]
Then
\begin{equation}
\Rcal^{\mathscr H}_\Phi(S)\subseteq\Rcal_\Phi(S).
\label{eq:app-restricted-exact-subfamily}
\end{equation}
Hence failure of the maximal factorization condition is an access obstruction shared by every downstream architecture using only $Q_K(X)$, while membership in $\Rcal_\Phi(S)$ need not imply membership in $\Rcal^{\mathscr H}_\Phi(S)$ for a restricted architecture class.

\subsection{State-relative presentation equivalence}

For a probability state $\mu$ on $X$, define
\[
(K,h)\equiv_\mu(K',h')
\quad\Longleftrightarrow\quad
hQ_K=h'Q_{K'}
\qquad\mu\text{-almost surely}.
\]

\begin{proposition}[State-relative presentation equivalence]
\label{prop:app-state-relative-presentation-equivalence}
For every probability state $\mu$, the relation $\equiv_\mu$ is an equivalence relation on receiver-level presentations with common action space $\mathsf A$. Moreover:
\begin{enumerate}[label=(\roman*)]
\item every refinement arrow in $\int\Dec_{\mathsf A}$ relates equivalent presentations pointwise and therefore under every state;
\item if $\mu\sim\nu$ have the same null sets, then $\equiv_\mu$ and $\equiv_\nu$ coincide;
\item if $\bar\mu=\sum_{c=1}^m\lambda_c\mu_c$ with every $\lambda_c>0$, then
\[
(K,h)\equiv_{\bar\mu}(K',h')
\quad\Longleftrightarrow\quad
(K,h)\equiv_{\mu_c}(K',h')
\text{ for every }c.
\]
\end{enumerate}
\end{proposition}

\begin{proof}
Reflexivity, symmetry, and transitivity are inherited from almost-sure equality of measurable maps. Part (i) follows from \eqref{eq:app-refinement-preserves-action}. For part (ii), the disagreement set
\[
B:=\{x:h(Q_K(x))\neq h'(Q_{K'}(x))\}
\]
is measurable and is $\mu$-null if and only if it is $\nu$-null. For part (iii),
\[
\bar\mu(B)=\sum_{c=1}^m\lambda_c\mu_c(B).
\]
With all coefficients strictly positive, this vanishes exactly when every $\mu_c(B)$ vanishes.
\end{proof}

These properties justify the main text's hierarchy
\[
K\longrightarrow(K,h)\longrightarrow[hQ_K]_\mu
\]
and make architecture-specific specialization explicit without requiring the categorical details to interrupt the main argument.

\section{Finite distribution geometry}
\label{app:finite-distribution-geometry}

This appendix records finite-space consequences of the distribution-relative receiver Bayes risk that are useful for structural selection analysis but not needed for the main theorem spine.

Assume throughout that $X$, $Y$, and all receiver images are finite and specialize to
\[
\mathsf A=Y,
\qquad
\ell(y,\hat y)=\mathbf 1[y\neq\hat y].
\]

\subsection{Threshold stability away from a boundary}

For a tolerance $\varepsilon\in[0,1]$, define
\[
\gamma_{\mu,\varepsilon}
:=
\min_{K\subseteq J}
|\delta_\mu(K)-\varepsilon|.
\]

\begin{corollary}[Localization-family stability radius]
\label{cor:app-localization-family-stability}
If $\gamma_{\mu,\varepsilon}>0$ and
\[
d_{\TV}(\mu,\nu)<\gamma_{\mu,\varepsilon},
\]
then
\[
\Rcal_{\mu,\varepsilon}
=
\Rcal_{\nu,\varepsilon}.
\]
Hence their inclusion-minimal localization antichains are also equal.
\end{corollary}

\begin{proof}
The main-text TV-Lipschitz theorem implies
\[
|\delta_\mu(K)-\delta_\nu(K)|
<\gamma_{\mu,\varepsilon}
\]
for every $K$. Every defect value therefore remains on the same side of the threshold $\varepsilon$.
\end{proof}

\subsection{Finite receiver risk is piecewise linear}

For fixed $K$, write
\[
p_{z,y}^\mu
:=
\mu\{x:Q_K(x)=z,\ \Phi(x)=y\}.
\]

\begin{proposition}[Finite receiver risk is ordinary Bayes classification error]
\label{thm:app-bayes-defect-formula}
For every receiver family $K$,
\[
\delta_\mu(K)
=
1-\sum_{z\in Z_K}\max_{y\in Y}p_{z,y}^\mu.
\]
In particular, $\mu\mapsto\delta_\mu(K)$ is concave and piecewise linear on the probability simplex $\Delta(X)$.
\end{proposition}

\begin{proof}
For a rule $h:Z_K\to Y$,
\[
R_\mu(K,h)
=
1-\sum_{z\in Z_K}p_{z,h(z)}^\mu.
\]
Each receiver value can be optimized independently, so the largest correct mass at $z$ is $\max_y p_{z,y}^\mu$.
\end{proof}

\subsection{Polyhedral selection geometry}

\begin{proposition}[Finite localization cells]
\label{thm:app-polyhedral-localization-cells}
Fix a finite receiver family $J$ and tolerance $\varepsilon\in[0,1]$. There is a finite polyhedral subdivision of $\Delta(X)$ such that on the relative interior of every cell:
\begin{enumerate}[label=(\roman*)]
\item the optimizer set $\Opt_\mu(K)$ is constant for every $K\subseteq J$;
\item the thresholded localization family $\Rcal_{\mu,\varepsilon}$ is constant;
\item the inclusion-minimal localization antichain
\[
\Mcal_{\mu,\varepsilon}:=\Min_{\subseteq}\Rcal_{\mu,\varepsilon}
\]
is constant.
\end{enumerate}
\end{proposition}

\begin{proof}
For fixed $K$ and $h$, the map
\[
\mu\longmapsto R_\mu(K,h)
\]
is affine. Partition the finite decoder set into classes whose members induce identical affine risk functionals. Distinct classes can exchange optimizer status only on nontrivial tie hyperplanes
\[
R_\mu(K,h)=R_\mu(K,h'),
\]
and threshold membership can change only on nontrivial hyperplanes
\[
R_\mu(K,h)=\varepsilon.
\]
The finite arrangement of these hyperplanes together with the faces of $\Delta(X)$ yields a finite polyhedral subdivision. On each relative interior the ordering and equality pattern of all distinct risk functionals, and their signs relative to $\varepsilon$, are fixed. The optimizing equivalence classes are therefore fixed; every decoder inside one such class enters or leaves the optimizer set together, so the full set $\Opt_\mu(K)$ is fixed as well. The thresholded family and its inclusion-minimal members are then fixed on the same cell.
\end{proof}

We use \emph{localization cell} for a cell of such a subdivision and reserve \emph{input regime} or \emph{context} for the scientific population being analyzed. The distinction avoids overloading the word ``regime.''

The main text uses only the consequence needed for interpretation: quantitative receiver risks vary continuously under TV shift, while optimizer identities and thresholded selections can change discretely at decision boundaries.

\section{Further structure of the exact support shadow}
\label{app:exact-shadow-details}

This appendix records exact-support structure that is useful for auditing and extending the support shadow but is not required for the main narrative.

\subsection{Support transport}

For finite discrete spaces define
\[
\supp\mu:=\{x:\mu(x)>0\}.
\]

\begin{proposition}[Finite support shadow]
\label{prop:app-support-shadow-functor}
If $f:(X,\mu)\to(Y,\nu)$ is a morphism in the finite discrete subcategory of $\int\Prob$, then
\[
\supp\nu=f(\supp\mu).
\]
Consequently support defines a functor from finite state-decorated spaces to finite sets by restricting each deterministic map to the support of its source state.
\end{proposition}

\begin{proof}
For $y\in Y$,
\[
\nu(y)
=(f_*\mu)(y)
=\sum_{x:f(x)=y}\mu(x),
\]
which is positive exactly when some $x\in\supp\mu$ satisfies $f(x)=y$.
\end{proof}

\subsection{Measure-class invariance of exact realization}

Let
\[
\Rcal^{\mathrm{a.s.}}_\Phi(\mu)
:=
\{K\subseteq J:\exists h:Z_K\to Y\text{ measurable with }\Phi=hQ_K\ \mu\text{-a.s.}\}.
\]

\begin{proposition}[Exact localization factors through measure class]
\label{thm:app-measure-class-invariance}
If $\mu\sim\nu$, then
\[
\Rcal^{\mathrm{a.s.}}_\Phi(\mu)
=
\Rcal^{\mathrm{a.s.}}_\Phi(\nu).
\]
\end{proposition}

\begin{proof}
For fixed $K$ and $h$, let
\[
B_{K,h}:=\{x:\Phi(x)\neq h(Q_K(x))\}.
\]
The decoder realizes $\Phi$ exactly under $\mu$ iff $B_{K,h}$ is $\mu$-null, which under $\mu\sim\nu$ is equivalent to being $\nu$-null. Existence of a realizing decoder is therefore the same under both laws.
\end{proof}

Together with Proposition~\ref{prop:app-measure-class-pushforward}, this gives a process-level coarse-graining: equivalent source states remain equivalent at every deterministic cut, and exact almost-sure localization cannot distinguish their different positive weights.

\subsection{Support restriction and truth-valued localization}

For finite $S\subseteq X$, write
\[
\Rcal_\Phi(S)
:=
\{K\subseteq J:\exists h:Q_K(S)\to Y\text{ with }\Phi|_S=hQ_K|_S\}.
\]

\begin{proposition}[Support restriction]
\label{thm:app-support-restriction-shadow}
If $S\subseteq T$, then
\[
\Rcal_\Phi(T)\subseteq\Rcal_\Phi(S).
\]
\end{proposition}

\begin{proof}
Any decoder that works on $T$ still works after restriction to $S$.
\end{proof}

Write
\[
\Supp(X):=(\mathcal P(X),\subseteq),
\qquad
\Two=\{0<1\}.
\]
Exact support localization can be represented as
\begin{equation}
\Loc_{\Phi,Q}:
\Supp(X)^{\mathrm{op}}\times\Bool_J
\longrightarrow\Two,
\label{eq:app-support-shadow-profunctor}
\end{equation}
with $\Loc_{\Phi,Q}(S,K)=1$ exactly when $K\in\Rcal_\Phi(S)$. Currying gives
\begin{equation}
\mathscr R_{\Phi,Q}:
\Supp(X)^{\mathrm{op}}
\longrightarrow
\Up(\Bool_J),
\qquad
S\longmapsto\Rcal_\Phi(S).
\label{eq:app-support-shadow-localization-functor}
\end{equation}
This truth-valued object remembers exact factorization on support and forgets all probability weights.

\subsection{Discernibility clutters and blockers}

For a map $f:S\to T$, write
\[
\Eq(f)
:=
\{(x,x')\in S^2:f(x)=f(x')\}.
\]
For a phenomenon-distinguished pair $x,x'\in S$, define
\[
E(x,x')
:=
\{j\in J:Q_j(x)\neq Q_j(x')\},
\]
and the distinction hypergraph
\[
\Hcal_\Phi(S)
:=
\{E(x,x'):x,x'\in S,\ \Phi(x)\neq\Phi(x')\}.
\]
For a hypergraph $\Hcal\subseteq2^J$, let
\[
\Hit(\Hcal)
:=
\{K\subseteq J:K\cap E\neq\varnothing\text{ for every }E\in\Hcal\}.
\]
The main text records the equivalence
\[
K\in\Rcal_\Phi(S)
\Longleftrightarrow
\Eq(Q_K|_S)\subseteq\Eq(\Phi|_S)
\Longleftrightarrow
K\text{ hits every }E\in\Hcal_\Phi(S).
\]

Deleting hyperedges that strictly contain another gives the clutter
\[
\Ccal_\Phi(S):=\Min_{\subseteq}\Hcal_\Phi(S).
\]
If
\[
b(\Ccal):=\Min_{\subseteq}\Hit(\Ccal)
\]
is the blocker, then the inclusion-minimal exact localizations satisfy
\begin{equation}
\Mcal_\Phi(S)=b(\Ccal_\Phi(S)).
\label{eq:app-minimal-localizations-blocker-shadow}
\end{equation}
This is classical transversal / decision-reduct structure. It makes explicit why an inclusion-minimal circuit need not be the unique minimal circuit.

\subsection{Uniform empirical defect and the classical \texorpdfstring{$g_3$}{g3} error}

Suppose $S=\{x_1,\ldots,x_n\}$ carries the uniform empirical law and use $0$--$1$ loss. For each receiver value $z$, let
\[
n_{z,y}
:=
\bigl|\{x\in S:Q_K(x)=z,\ \Phi(x)=y\}\bigr|.
\]
The finite Bayes formula gives
\begin{equation}
\delta_{\mathrm{unif}(S)}(K)
=
1-\frac1n\sum_z\max_y n_{z,y}.
\label{eq:app-empirical-defect-g3}
\end{equation}
To make the functional dependency $K\to d$ exact by deleting rows, one may retain at most one decision value inside each $Q_K$-equivalence class. The largest exact subtable therefore keeps $\sum_z\max_y n_{z,y}$ rows. Hence \eqref{eq:app-empirical-defect-g3} is exactly the classical $g_3$ deletion error for an approximate functional dependency \citep{kivinen1995approximate}.

Variable-precision and decision-theoretic rough-set models likewise study probabilistic or loss-sensitive relaxations of exact reducts \citep{ziarko1993variable,zhao2008attribute}. The paper's distributional content lies upstream of this finite empirical shadow: the probability state is transported through one neural process and receiver refinement is compared with context invariance before these coarse-grainings are taken.

\section{Further fixed-network universality consequences}
\label{app:fixed-or-extensions}

This appendix records consequences of the fixed-OR support-universality theorem that are stronger than needed for the main narrative.

For a support $S$ of the fixed OR network, define its \emph{exact localization order} by
\begin{equation}
\lambda_\Phi(S)
:=
\min\{|K|:K\in\Rcal_\Phi(S)\}.
\label{eq:app-exact-localization-order}
\end{equation}
The family $\Rcal_\Phi(S)$ is nonempty because the full coordinate set $J$ always determines the OR output on $S$.

\begin{corollary}[Arbitrary monotone localization-order profile]
\label{cor:app-arbitrary-order-profile-shadow}
Let $P$ be finite and let
\[
d:P\to\{0,\ldots,|J|\}
\]
satisfy
\[
p\le q\Longrightarrow d(p)\le d(q).
\]
Then the fixed network
\[
\NN_J:
\{0,1\}^J
\xrightarrow{\Id}
\{0,1\}^J
\xrightarrow{\mathrm{OR}}
\{0,1\}
\]
admits supports $S_p$ such that
\begin{equation}
p\le q
\Longrightarrow
S_p\subseteq S_q,
\qquad
\lambda_\Phi(S_p)=d(p).
\label{eq:app-monotone-support-profile}
\end{equation}
\end{corollary}

\begin{proof}
Set
\[
F_p:=\{K\subseteq J:|K|\ge d(p)\}
\]
and
\[
S_p:=\sigma_{\mathrm{OR}}(F_p).
\]
Each $F_p$ is nonempty and upward closed. If $p\le q$, then $d(p)\le d(q)$ and hence
\[
F_q\subseteq F_p.
\]
Every hitting set of $F_p$ therefore hits $F_q$, so
\[
\Hit(F_p)\subseteq\Hit(F_q).
\]
By the definition of $\sigma_{\mathrm{OR}}$, this gives
\[
S_p\subseteq S_q.
\]
Theorem~\ref{thm:fixed-or-support-universality} gives $\Rcal_\Phi(S_p)=F_p$, whose smallest members have cardinality exactly $d(p)$. Thus $\lambda_\Phi(S_p)=d(p)$.
\end{proof}

The result shows that the fixed OR network can realize not only arbitrary admissible localization families at one support, but arbitrary finite monotone profiles of minimum exact receiver-set size along nested support diagrams. The main text omits this strengthening so that the fixed-OR theorem itself remains the final major theorem and its direct mechanistic consequence is visually dominant.

\section{Access-image descent for exact mechanisms}
\label{app:access-image-descent}

This appendix preserves the local-to-global exact compatibility theory. It is mathematically downstream of the main support-shadow story and is therefore removed from the main theorem spine, but no result is discarded.

Suppose every context individually admits an exact receiver-level decoder. It is still possible that no single decoder works across their union. The obstruction does not have to appear at the source: a noninjective internal access can identify states from different contexts that were distinct before the access. Local exact descriptions must therefore be compared on the overlaps created by the receiver itself.

\subsection{Receiver access can create new cross-context collisions}

For a map $f:X\to Y$ in $\Set$, direct image and inverse image define
\[
\Sigma_f(A):=f(A),
\qquad
f^*(B):=f^{-1}(B),
\]
with
\begin{equation}
\Sigma_f\dashv f^*.
\label{eq:app-direct-inverse-adjunction}
\end{equation}
Direct image preserves unions but only satisfies
\begin{equation}
f(A\cap B)
\subseteq
f(A)\cap f(B).
\label{eq:app-direct-image-created-overlap}
\end{equation}
The inclusion can be strict exactly because $f$ may identify a point of $A\setminus B$ with a point of $B\setminus A$. Thus two disjoint or weakly overlapping source regimes can acquire a larger overlap after a network prefix or receiver access.

For mechanistic compatibility, the relevant overlap is therefore
\[
Q_K(S_i)\cap Q_K(S_j),
\]
not merely $S_i\cap S_j$.

\subsection{Exact context compatibility is descent on the access-image cover}

Let
\[
S=\bigcup_{i\in I}S_i
\]
be a finite cover and let
\[
L:S\twoheadrightarrow U=L(S)
\]
be an access. Put
\[
U_i:=L(S_i).
\]
Assume $L|_{S_i}$ exactly realizes $\Phi|_{S_i}$ for every $i$. Each context then has a unique local decoder
\[
h_i:U_i\to Y
\]
with
\[
\Phi|_{S_i}=h_iL|_{S_i}.
\]
The local decoder domains form a cover
\[
U=\bigcup_iU_i.
\]

\begin{theorem}[Exact context compatibility is an access-image global-section condition]
\label{thm:app-access-image-descent}
Under the preceding assumptions, $L$ exactly realizes $\Phi$ on $S$ if and only if
\begin{equation}
 h_i|_{U_i\cap U_j}
=
h_j|_{U_i\cap U_j}
\qquad\text{for all }i,j.
\label{eq:app-access-image-descent}
\end{equation}
When a global decoder exists, it is unique on $U$.
\end{theorem}

\begin{proof}
If $h:U\to Y$ is a global decoder, uniqueness of factorization through the surjections $L|_{S_i}:S_i\twoheadrightarrow U_i$ gives $h_i=h|_{U_i}$, so compatibility is necessary.

Conversely, if the local decoders agree on every overlap, define
\[
h(u):=h_i(u)
\]
for any $i$ with $u\in U_i$. The overlap condition makes this well defined, and for $x\in S_i$,
\[
h(L(x))=h_i(L(x))=\Phi(x).
\]
Surjectivity gives uniqueness on $U$.
\end{proof}

The theorem is ordinary function gluing. The mechanistic identification is the domain on which it must be applied: the cover is generated by access images, not by source contexts before the access.

\subsection{Sheaf formulation}

On the poset of subsets of $U$, define
\[
\mathscr D_Y(V):=\{h:V\to Y\}
\]
with restriction along inclusions. For the ordinary union coverage, $\mathscr D_Y$ is the usual sheaf of $Y$-valued functions. The local decoders are sections
\[
h_i\in\mathscr D_Y(U_i),
\]
and Theorem~\ref{thm:app-access-image-descent} says that exact global realization is equivalent to their extension to one section of $\mathscr D_Y(U)$.

This has the same local-to-global form used in sheaf-theoretic contextuality \citep{abramsky2011sheaf}, but the objects are different. Here the cover is generated by access images of input regimes through one deterministic receiver interface, and the local data are decoder functions constrained by one phenomenon. No measurement scenario, no-signalling condition, or hidden-variable semantics is assumed. The shared mathematics is the global-section obstruction; the mechanistic content is the access-induced cover.

\subsection{Descent as the zero-loss endpoint of context value}

Assume finite spaces with
\[
\mathsf A=Y,
\qquad
\ell(y,\hat y)=\mathbf1[y\neq\hat y].
\]
Let $\mu_i$ have full support $S_i$ and suppose
\[
\delta_{\mu_i}(K)=0
\]
for every context. Let
\[
\bar\mu
=
\sum_i\lambda_i\mu_i,
\qquad
\lambda_i>0,
\quad
\sum_i\lambda_i=1.
\]
Then $\supp\bar\mu=\bigcup_iS_i$.

\begin{corollary}[Zero context value becomes exact descent]
\label{cor:app-descent-zero-contextuality}
For a fixed receiver family $K$, the following are equivalent:
\begin{enumerate}[label=(\roman*)]
\item $\delta_{\bar\mu}(K)=0$;
\item the contexts admit one common zero-risk decoder through $Q_K$;
\item the local exact decoders agree on all access-image intersections $Q_K(S_i)\cap Q_K(S_j)$;
\item the local decoder sections extend to a global section of $\mathscr D_Y$ on $Q_K(\bigcup_iS_i)$;
\item $Q_K$ exactly realizes $\Phi$ on $\bigcup_iS_i$.
\end{enumerate}
\end{corollary}

\begin{proof}
Finite zero defect is equivalent to exact realization on support. Since every local optimum is zero and every mixture coefficient is positive, a zero-risk mixture decoder is exactly a rule with zero risk in every context, giving (i)$\Leftrightarrow$(ii). Theorem~\ref{thm:app-access-image-descent} gives the equivalence of (ii), (iii), and (v), while the sheaf condition gives (iii)$\Leftrightarrow$(iv).
\end{proof}

For positive loss, $\Gamma_K$ measures how much Bayes risk is saved by revealing context. At zero loss, that quantitative question collapses to the exact extension problem above. Thus context compatibility has a local-to-global endpoint even though the main paper no longer spends a full section developing it.


\end{document}
